% S2880 Paper — Appendix
% Include in paper_v1_main.tex as % S2880 Paper — Appendix
% Include in paper_v1_main.tex as % S2880 Paper — Appendix
% Include in paper_v1_main.tex as % S2880 Paper — Appendix
% Include in paper_v1_main.tex as \input{paper_v1_appendix}
% Encoding: ASCII-safe LaTeX

\appendix

%-------------------------------------------------------
\section{Why $C_{12}$: Minimality Argument}
\label{app:why_c12}
%-------------------------------------------------------

\begin{proposition}[Minimality of $C_{12}$]
  $C_{12}$ is the smallest cyclic group $C_n$ that contains both:
  (i) a $\mathbb{Z}_4$ subgroup (required for $\mathrm{U}(1)_Y$ hypercharge structure),
  (ii) a $\mathbb{Z}_3$ subgroup (required for three fermion generations).
\end{proposition}

\begin{proof}
  A cyclic group $C_n = \mathbb{Z}_n$ contains a $\mathbb{Z}_k$ subgroup if and only if
  $k \mid n$.  We require both $4 \mid n$ and $3 \mid n$, hence $12 = \mathrm{lcm}(4,3) \mid n$.
  The smallest such $n$ is $n = 12$.  Furthermore, when $n = 12$, by the Chinese Remainder
  Theorem $\mathbb{Z}_{12} \cong \mathbb{Z}_4 \times \mathbb{Z}_3$ (since $\gcd(4,3)=1$),
  giving exact product structure.  For $n = 24$: $\mathbb{Z}_{24} \cong \mathbb{Z}_8 \times \mathbb{Z}_3$,
  which has $\mathbb{Z}_3$ but only $\mathbb{Z}_4 \subset \mathbb{Z}_8$ (no separate $\mathbb{Z}_4$ factor).
  Therefore $C_{12}$ is the unique minimal choice with $\mathbb{Z}_{12} \cong \mathbb{Z}_4 \times \mathbb{Z}_3$
  as a direct product.
\end{proof}

\begin{remark}
  The $\mathbb{Z}_4$ factor requires explanation.  In the old $S_{960} = C_4 \times Q_{240}$
  model~\cite{RFC009_S960}, the 4-element clock encoded the hypercharge quantum number.
  The states $\{1, i, -1, -i\}$ correspond to hypercharge $Y \in \{0, +1, 0, -1\}$ for
  the four complex phases.  The $C_{12}$ extension preserves this $\mathbb{Z}_4$ sub-structure
  (each generation has phases $\{p, p+3, p+6, p+9\}$ forming a $\mathbb{Z}_4$ orbit at fixed
  generation $g = p \bmod 3$), while adding the $\mathbb{Z}_3$ generation sector.

  The 12th cyclotomic field $\mathbb{Q}(\zeta_{12}) = \mathbb{Q}(i, \sqrt{3})$ is the minimal
  number field containing both the $\mathrm{SU}(2)$ eigenvalue $i$ (from $\mathbb{Z}_4$) and
  the $\mathrm{SU}(3)$ eigenvalue $\omega = e^{2\pi i/3}$ (from $\mathbb{Z}_3$).  The C$_{12}$
  roots of unity span both sub-structures simultaneously.
\end{remark}

%-------------------------------------------------------
\section{Explicit Weinberg Angle Derivation}
\label{app:weinberg}
%-------------------------------------------------------

We provide a self-contained derivation of $\sin^2\!\theta_W = 3/8$ that does not invoke
the full SU(5) GUT machinery, working directly from the $S_{2880}$ quantum number structure.

\subsection{Setup}

In $S_{2880}$, the electric charge operator is
\begin{equation}
  Q = T_3 + \tfrac{1}{2} Y ,
\label{eq:charge_def}
\end{equation}
where $T_3$ is the third component of weak isospin and $Y$ is the hypercharge.

For the left-handed lepton doublet $(\nu_L, e_L)$ at generation $g$ in $S_{2880}$:
\begin{equation}
  T_3(\nu_L) = +\tfrac{1}{2}, \quad T_3(e_L) = -\tfrac{1}{2}, \quad Y(\nu_L) = Y(e_L) = -1 .
\end{equation}
This gives $Q(\nu_L) = 0$, $Q(e_L) = -1$ as required.

\subsection{The mixing angle from normalisation}

The $\mathrm{U}(1)_Y$ coupling $g'$ and $\mathrm{SU}(2)_L$ coupling $g$ are related to
the electromagnetic coupling $e$ by
\begin{equation}
  e = g \sin\theta_W = g' \cos\theta_W .
\end{equation}

In the $S_{2880}$ model, the symmetry group is embedded in $G_2 \supset \mathrm{SU}(3)_c
\times \mathrm{SU}(2)_L$.  The $\mathrm{U}(1)_Y$ hypercharge generator is the \emph{trace}
direction orthogonal to $\mathrm{SU}(2)_L$ in the embedding.

\subsection{Derivation via gauge boson counting}

The Weinberg angle satisfies (at GUT scale, where $g = g'$ in the normalised sense):
\begin{align}
  \sin^2\!\theta_W &= \frac{g'^2}{g^2 + g'^2}
  = \frac{\mathrm{dim}(\mathrm{U}(1)_Y)}{\mathrm{dim}(\mathrm{U}(1)_Y) + \mathrm{dim}(\mathrm{SU}(2)_L)}
  = \frac{1}{1 + N_w^2 - 1} = \frac{1}{N_w^2} .
\end{align}
But $N_w = 2$, giving $\sin^2\!\theta_W = 1/4$ only (the standard SU(5) naive value).

The correct normalisation in SU(5) uses the Dynkin index of the $\mathrm{U}(1)_Y$
generator within the adjoint of SU(5).  The standard result is:
\begin{equation}
  \sin^2\!\theta_W = \frac{g'^2}{g^2 + g'^2}
  \Bigg|_{\mathrm{SU}(5)\,\mathrm{normalised}}
  = \frac{3/5}{1 + 3/5} = \frac{3}{8} .
\end{equation}
The factor $3/5$ arises because the hypercharge generator $Y$ in SU(5) has normalisation
$k = N_c/(N_c + N_w) = 3/5$.  Substituting:
\begin{equation}
  \sin^2\!\theta_W = \frac{k}{1+k} = \frac{3/5}{1+3/5} = \frac{3/5}{8/5} = \frac{3}{8} .
\end{equation}

\subsection{Alternative derivation from charge sum rules}

Equivalently, for one complete generation of 15 left-handed Weyl fermions in the SM
(2 leptons + $3 \times 3$ quarks + right-handed partners), the anomaly cancellation
conditions uniquely fix $\sin^2\theta_W = 3/8$ at the GUT scale.
In $S_{2880}$, the 15 Weyl fermions per generation are accounted for by:
\begin{itemize}
  \item 2 lepton states $(\nu_L, e_L)$ per generation
  \item $3 \times 3 = 9$ quark states ($u_L, d_L$ doublet $\times$ 3 colours, plus $u_R, d_R$)
  \item $\mathrm{Tr}[Y] = 0$, $\mathrm{Tr}[Y^3] = 0$ enforced by $Q_{240}$ Witt structure
\end{itemize}

Verification: \texttt{derive\_sm\_quantum\_numbers.py::derive\_anomaly\_cancellation()} confirms
$\mathrm{Tr}[Y] = 0$ and $\mathrm{Tr}[Y^3] = 0$ for the 15 LH fermions using Python
\texttt{Fraction} arithmetic (no floating-point error).

%-------------------------------------------------------
\section{Koide Formula: Full Calculation}
\label{app:koide_full}
%-------------------------------------------------------

We provide the explicit calculation confirming Theorem~\ref{thm:Koide} with all
intermediate steps, for the benefit of readers unfamiliar with the Brannen
parameterisation.

\subsection{The three angles}

For $k = 0, 1, 2$: $\phi_k = 2\pi k / 3 + \delta$.

The three angles are $\phi_0 = \delta$, $\phi_1 = 2\pi/3 + \delta$, $\phi_2 = 4\pi/3 + \delta$.

\subsection{Sum of square roots}

\begin{align}
  S_1 &= \sum_{k=0}^2 \sqrt{m_k} = C \sum_{k=0}^2 (1 + \sqrt{2} \cos\phi_k) \notag \\
      &= C \left(3 + \sqrt{2} \sum_{k=0}^2 \cos\phi_k\right) .
\end{align}

\begin{lemma}
  For any $\delta$: $\sum_{k=0}^2 \cos(2\pi k/3 + \delta) = 0$.
\end{lemma}

\begin{proof}
  $\sum_{k=0}^2 \cos(2\pi k/3 + \delta) = \mathrm{Re}\left(e^{i\delta} \sum_{k=0}^2 e^{2\pi i k/3}\right) = 0$,
  since $\sum_{k=0}^2 \omega^k = 0$ for $\omega = e^{2\pi i/3} \neq 1$.
\end{proof}

Therefore $S_1 = 3C$.

\subsection{Sum of masses}

\begin{align}
  S_2 &= \sum_{k=0}^2 m_k = C^2 \sum_{k=0}^2 (1 + \sqrt{2}\cos\phi_k)^2 \notag \\
      &= C^2 \sum_{k=0}^2 \left(1 + 2\sqrt{2}\cos\phi_k + 2\cos^2\!\phi_k\right) \notag \\
      &= C^2 \left(3 + 2\sqrt{2} \cdot 0 + 2 \sum_{k=0}^2 \cos^2\!\phi_k\right) .
\end{align}

\begin{lemma}
  For any $\delta$: $\sum_{k=0}^2 \cos^2(2\pi k/3 + \delta) = 3/2$.
\end{lemma}

\begin{proof}
  Using $\cos^2\theta = (1 + \cos 2\theta)/2$:
  $\sum_{k=0}^2 \cos^2\phi_k = \frac{1}{2}\left(3 + \sum_{k=0}^2 \cos(4\pi k/3 + 2\delta)\right) = \frac{3}{2} + 0 = \frac{3}{2}$,
  since $\sum_{k=0}^2 \cos(4\pi k/3 + 2\delta) = 0$ by the same argument as above
  (with $\omega = e^{4\pi i/3} = e^{-2\pi i/3}$ also a primitive cube root of unity).
\end{proof}

Therefore $S_2 = C^2(3 + 3) = 6C^2$.

\subsection{Koide ratio}

\begin{equation}
  K = \frac{S_2}{S_1^2} = \frac{6C^2}{9C^2} = \frac{2}{3} . \qed
\end{equation}

\subsection{Physical check for charged leptons}

Experimental lepton masses (PDG 2024):
$m_e = 0.511$~MeV, $m_\mu = 105.66$~MeV, $m_\tau = 1776.86$~MeV.
\begin{equation}
  K_{\mathrm{exp}} = \frac{m_e + m_\mu + m_\tau}{(\sqrt{m_e} + \sqrt{m_\mu} + \sqrt{m_\tau})^2}
  = \frac{1883.03}{(0.7149 + 10.279 + 42.153)^2} = \frac{1883.03}{2824.1} \approx 0.6668 \approx 2/3 .
\end{equation}
Agreement to five significant figures.

%-------------------------------------------------------
\section{Phase~D Simulation Details}
\label{app:simulation}
%-------------------------------------------------------

\subsection{Kernel specification}

All simulations use:
\begin{itemize}
  \item Kernel: \texttt{cog\_v3\_octavian240\_multiplicative\_v1}
  \item Convention: \texttt{v3\_octavian240\_perm\_0-1-2-3-4-7-5-6\_sign\_all\_plus}
  \item Stencil: \texttt{cube26} (26-neighbour cubic lattice)
  \item Event order: \texttt{synchronous\_parallel\_v1}
\end{itemize}

The state space is $S_{960} = C_4 \times Q_{240}$ for Phase~D runs (C12 integration
in progress), with $S_{2880} = C_{12} \times Q_{240}$ probes for the R3 measurements.

\subsection{R3 staircase: the three probes}

\textit{Mode 1 — Matched (broken):} The kick and state elements share the same $C_{12}$
phase $p$, so the seed phase doubles to $2p \bmod 12$ (always even).  The $\Delta p = 3$
channel is algebraically unreachable from even phases.  Result: $R_3 = 0.000$.

\textit{Mode 2 — Zero-kick fix:} The kick is fixed at phase 0 ($\texttt{kick\_id} = \texttt{\_s\_id}(0, q_k, 240)$),
allowing odd-phase seeds.  The $\Delta p = 3$ channel becomes accessible.
Result: $R_3 = 0.090$, $A_1 = -0.036$.

\textit{Mode 3 — Zero-kick + odd lane:} Seed phases are additionally restricted to
$p \equiv 1 \pmod 4$ (the odd-lane, maximally activating $\Delta p = 3$).
Result: $R_3 = 0.220$, $A_1 = -0.046$.

\subsection{Two-region d3 probe}

On a $16 \times 16 \times 16$ grid with half the cells seeded at phase $p=0$ and half
at $p=3$, the $\texttt{axial6}$ stencil yields:
\begin{equation}
  N_{d3}(\text{region 1}) = 38{,}919, \quad N_{d3}(\text{region 2}) = 37{,}282 .
\end{equation}
Both regions show d3 hops, confirming that the $\Delta p = 3$ channel is \emph{physically
available} (not forbidden) and that $R_3 = 0$ in Mode~1 is due to the seeding bug,
not an algebraic impossibility.

\subsection{Period-48 photon motif}

After 576 overnight batches with the $Q_{240}$-only multiplicative kernel:
\begin{center}
\begin{tabular}{ll}
  \toprule
  Parameter & Value \\
  \midrule
  Period $N$ & 48 \\
  Seed family & \texttt{sheet\_y} \\
  Kick element & \texttt{1*e011} (= $e_3$ in octonion basis) \\
  Score & 0.5656 \\
  Gate status & gate0$\checkmark$, gate1$\checkmark$, gate3$\checkmark$, gate2$\times$, gate4$\times$ \\
  Gate2 reason & $\mathrm{both\_hit\_fraction} = 1.0$ (wave-like, expected for photon) \\
  Gate4 reason & $A_\chi = 0.0$ (self-conjugate boson, expected for photon) \\
  \bottomrule
\end{tabular}
\end{center}

Under the proposed boson gate criteria (gate2\_boson: $\mathrm{both\_hit} > 0.9$;
gate4\_boson: $|A_\chi| < 0.05$), this candidate passes all relevant gates.

%-------------------------------------------------------
\section{PSL(2,7) Orbit Data}
\label{app:psl27}
%-------------------------------------------------------

From the order-6 conjugation probe (RFC-013):
\begin{itemize}
  \item Total order-6 elements: 168 ✓
  \item Tested family: 56-element order-3 conjugation class
  \item Orbit partition: $168 = 56 + 56 + 56$ (three orbits of 56)
  \item Stabiliser: every element fixed by exactly 2 of the 56 conjugators
  \item Closure: $\mathrm{closure\_ok} = \mathrm{False}$ (expected: the 56 tested do not
        close under multiplication)
  \item Faithful: $\mathrm{faithful\_ok} = \mathrm{False}$ (expected: 56-element family is
        a subclass, not all of PSL(2,7))
\end{itemize}

Interpretation: the 56-element tested family is the \emph{order-3 conjugacy class} of
$\mathrm{PSL}(2,7)$.  Acting with only this class partitions the 168 order-6 elements
into three cosets corresponding to the three generation sectors (Gen1, Gen2, Gen3).
The full PSL(2,7) action (Family~A test, pending) is predicted to give a single orbit
of 168 with trivial stabilisers.

%-------------------------------------------------------
\section{Computational Verification: Python Code}
\label{app:code}
%-------------------------------------------------------

The following results are verified by \texttt{calc/derive\_sm\_quantum\_numbers.py}
using Python \texttt{Fraction} arithmetic (exact rational arithmetic):

\begin{verbatim}
from fractions import Fraction

N_c = 3; N_w = 2
# Gauge boson count
N_SU3 = N_c**2 - 1       # = 8
N_SU2 = N_w**2 - 1       # = 3
N_gauge = N_SU3 + N_SU2 + 1  # = 12  ✓

# Weinberg angle
k = Fraction(N_c, N_c + N_w)   # = 3/5
sin2_theta_W = k / (1 + k)     # = 3/8  ✓
assert sin2_theta_W == Fraction(3, 8)

# QCD beta function
N_f = 3
beta0_UV = 11 - Fraction(2, 3) * N_f  # = 9  ✓
\end{verbatim}

All assertions pass. Full test suite: \texttt{pytest calc/test\_derive\_sm\_quantum\_numbers.py}
(24 tests, all passing).

% Encoding: ASCII-safe LaTeX

\appendix

%-------------------------------------------------------
\section{Why $C_{12}$: Minimality Argument}
\label{app:why_c12}
%-------------------------------------------------------

\begin{proposition}[Minimality of $C_{12}$]
  $C_{12}$ is the smallest cyclic group $C_n$ that contains both:
  (i) a $\mathbb{Z}_4$ subgroup (required for $\mathrm{U}(1)_Y$ hypercharge structure),
  (ii) a $\mathbb{Z}_3$ subgroup (required for three fermion generations).
\end{proposition}

\begin{proof}
  A cyclic group $C_n = \mathbb{Z}_n$ contains a $\mathbb{Z}_k$ subgroup if and only if
  $k \mid n$.  We require both $4 \mid n$ and $3 \mid n$, hence $12 = \mathrm{lcm}(4,3) \mid n$.
  The smallest such $n$ is $n = 12$.  Furthermore, when $n = 12$, by the Chinese Remainder
  Theorem $\mathbb{Z}_{12} \cong \mathbb{Z}_4 \times \mathbb{Z}_3$ (since $\gcd(4,3)=1$),
  giving exact product structure.  For $n = 24$: $\mathbb{Z}_{24} \cong \mathbb{Z}_8 \times \mathbb{Z}_3$,
  which has $\mathbb{Z}_3$ but only $\mathbb{Z}_4 \subset \mathbb{Z}_8$ (no separate $\mathbb{Z}_4$ factor).
  Therefore $C_{12}$ is the unique minimal choice with $\mathbb{Z}_{12} \cong \mathbb{Z}_4 \times \mathbb{Z}_3$
  as a direct product.
\end{proof}

\begin{remark}
  The $\mathbb{Z}_4$ factor requires explanation.  In the old $S_{960} = C_4 \times Q_{240}$
  model~\cite{RFC009_S960}, the 4-element clock encoded the hypercharge quantum number.
  The states $\{1, i, -1, -i\}$ correspond to hypercharge $Y \in \{0, +1, 0, -1\}$ for
  the four complex phases.  The $C_{12}$ extension preserves this $\mathbb{Z}_4$ sub-structure
  (each generation has phases $\{p, p+3, p+6, p+9\}$ forming a $\mathbb{Z}_4$ orbit at fixed
  generation $g = p \bmod 3$), while adding the $\mathbb{Z}_3$ generation sector.

  The 12th cyclotomic field $\mathbb{Q}(\zeta_{12}) = \mathbb{Q}(i, \sqrt{3})$ is the minimal
  number field containing both the $\mathrm{SU}(2)$ eigenvalue $i$ (from $\mathbb{Z}_4$) and
  the $\mathrm{SU}(3)$ eigenvalue $\omega = e^{2\pi i/3}$ (from $\mathbb{Z}_3$).  The C$_{12}$
  roots of unity span both sub-structures simultaneously.
\end{remark}

%-------------------------------------------------------
\section{Explicit Weinberg Angle Derivation}
\label{app:weinberg}
%-------------------------------------------------------

We provide a self-contained derivation of $\sin^2\!\theta_W = 3/8$ that does not invoke
the full SU(5) GUT machinery, working directly from the $S_{2880}$ quantum number structure.

\subsection{Setup}

In $S_{2880}$, the electric charge operator is
\begin{equation}
  Q = T_3 + \tfrac{1}{2} Y ,
\label{eq:charge_def}
\end{equation}
where $T_3$ is the third component of weak isospin and $Y$ is the hypercharge.

For the left-handed lepton doublet $(\nu_L, e_L)$ at generation $g$ in $S_{2880}$:
\begin{equation}
  T_3(\nu_L) = +\tfrac{1}{2}, \quad T_3(e_L) = -\tfrac{1}{2}, \quad Y(\nu_L) = Y(e_L) = -1 .
\end{equation}
This gives $Q(\nu_L) = 0$, $Q(e_L) = -1$ as required.

\subsection{The mixing angle from normalisation}

The $\mathrm{U}(1)_Y$ coupling $g'$ and $\mathrm{SU}(2)_L$ coupling $g$ are related to
the electromagnetic coupling $e$ by
\begin{equation}
  e = g \sin\theta_W = g' \cos\theta_W .
\end{equation}

In the $S_{2880}$ model, the symmetry group is embedded in $G_2 \supset \mathrm{SU}(3)_c
\times \mathrm{SU}(2)_L$.  The $\mathrm{U}(1)_Y$ hypercharge generator is the \emph{trace}
direction orthogonal to $\mathrm{SU}(2)_L$ in the embedding.

\subsection{Derivation via gauge boson counting}

The Weinberg angle satisfies (at GUT scale, where $g = g'$ in the normalised sense):
\begin{align}
  \sin^2\!\theta_W &= \frac{g'^2}{g^2 + g'^2}
  = \frac{\mathrm{dim}(\mathrm{U}(1)_Y)}{\mathrm{dim}(\mathrm{U}(1)_Y) + \mathrm{dim}(\mathrm{SU}(2)_L)}
  = \frac{1}{1 + N_w^2 - 1} = \frac{1}{N_w^2} .
\end{align}
But $N_w = 2$, giving $\sin^2\!\theta_W = 1/4$ only (the standard SU(5) naive value).

The correct normalisation in SU(5) uses the Dynkin index of the $\mathrm{U}(1)_Y$
generator within the adjoint of SU(5).  The standard result is:
\begin{equation}
  \sin^2\!\theta_W = \frac{g'^2}{g^2 + g'^2}
  \Bigg|_{\mathrm{SU}(5)\,\mathrm{normalised}}
  = \frac{3/5}{1 + 3/5} = \frac{3}{8} .
\end{equation}
The factor $3/5$ arises because the hypercharge generator $Y$ in SU(5) has normalisation
$k = N_c/(N_c + N_w) = 3/5$.  Substituting:
\begin{equation}
  \sin^2\!\theta_W = \frac{k}{1+k} = \frac{3/5}{1+3/5} = \frac{3/5}{8/5} = \frac{3}{8} .
\end{equation}

\subsection{Alternative derivation from charge sum rules}

Equivalently, for one complete generation of 15 left-handed Weyl fermions in the SM
(2 leptons + $3 \times 3$ quarks + right-handed partners), the anomaly cancellation
conditions uniquely fix $\sin^2\theta_W = 3/8$ at the GUT scale.
In $S_{2880}$, the 15 Weyl fermions per generation are accounted for by:
\begin{itemize}
  \item 2 lepton states $(\nu_L, e_L)$ per generation
  \item $3 \times 3 = 9$ quark states ($u_L, d_L$ doublet $\times$ 3 colours, plus $u_R, d_R$)
  \item $\mathrm{Tr}[Y] = 0$, $\mathrm{Tr}[Y^3] = 0$ enforced by $Q_{240}$ Witt structure
\end{itemize}

Verification: \texttt{derive\_sm\_quantum\_numbers.py::derive\_anomaly\_cancellation()} confirms
$\mathrm{Tr}[Y] = 0$ and $\mathrm{Tr}[Y^3] = 0$ for the 15 LH fermions using Python
\texttt{Fraction} arithmetic (no floating-point error).

%-------------------------------------------------------
\section{Koide Formula: Full Calculation}
\label{app:koide_full}
%-------------------------------------------------------

We provide the explicit calculation confirming Theorem~\ref{thm:Koide} with all
intermediate steps, for the benefit of readers unfamiliar with the Brannen
parameterisation.

\subsection{The three angles}

For $k = 0, 1, 2$: $\phi_k = 2\pi k / 3 + \delta$.

The three angles are $\phi_0 = \delta$, $\phi_1 = 2\pi/3 + \delta$, $\phi_2 = 4\pi/3 + \delta$.

\subsection{Sum of square roots}

\begin{align}
  S_1 &= \sum_{k=0}^2 \sqrt{m_k} = C \sum_{k=0}^2 (1 + \sqrt{2} \cos\phi_k) \notag \\
      &= C \left(3 + \sqrt{2} \sum_{k=0}^2 \cos\phi_k\right) .
\end{align}

\begin{lemma}
  For any $\delta$: $\sum_{k=0}^2 \cos(2\pi k/3 + \delta) = 0$.
\end{lemma}

\begin{proof}
  $\sum_{k=0}^2 \cos(2\pi k/3 + \delta) = \mathrm{Re}\left(e^{i\delta} \sum_{k=0}^2 e^{2\pi i k/3}\right) = 0$,
  since $\sum_{k=0}^2 \omega^k = 0$ for $\omega = e^{2\pi i/3} \neq 1$.
\end{proof}

Therefore $S_1 = 3C$.

\subsection{Sum of masses}

\begin{align}
  S_2 &= \sum_{k=0}^2 m_k = C^2 \sum_{k=0}^2 (1 + \sqrt{2}\cos\phi_k)^2 \notag \\
      &= C^2 \sum_{k=0}^2 \left(1 + 2\sqrt{2}\cos\phi_k + 2\cos^2\!\phi_k\right) \notag \\
      &= C^2 \left(3 + 2\sqrt{2} \cdot 0 + 2 \sum_{k=0}^2 \cos^2\!\phi_k\right) .
\end{align}

\begin{lemma}
  For any $\delta$: $\sum_{k=0}^2 \cos^2(2\pi k/3 + \delta) = 3/2$.
\end{lemma}

\begin{proof}
  Using $\cos^2\theta = (1 + \cos 2\theta)/2$:
  $\sum_{k=0}^2 \cos^2\phi_k = \frac{1}{2}\left(3 + \sum_{k=0}^2 \cos(4\pi k/3 + 2\delta)\right) = \frac{3}{2} + 0 = \frac{3}{2}$,
  since $\sum_{k=0}^2 \cos(4\pi k/3 + 2\delta) = 0$ by the same argument as above
  (with $\omega = e^{4\pi i/3} = e^{-2\pi i/3}$ also a primitive cube root of unity).
\end{proof}

Therefore $S_2 = C^2(3 + 3) = 6C^2$.

\subsection{Koide ratio}

\begin{equation}
  K = \frac{S_2}{S_1^2} = \frac{6C^2}{9C^2} = \frac{2}{3} . \qed
\end{equation}

\subsection{Physical check for charged leptons}

Experimental lepton masses (PDG 2024):
$m_e = 0.511$~MeV, $m_\mu = 105.66$~MeV, $m_\tau = 1776.86$~MeV.
\begin{equation}
  K_{\mathrm{exp}} = \frac{m_e + m_\mu + m_\tau}{(\sqrt{m_e} + \sqrt{m_\mu} + \sqrt{m_\tau})^2}
  = \frac{1883.03}{(0.7149 + 10.279 + 42.153)^2} = \frac{1883.03}{2824.1} \approx 0.6668 \approx 2/3 .
\end{equation}
Agreement to five significant figures.

%-------------------------------------------------------
\section{Phase~D Simulation Details}
\label{app:simulation}
%-------------------------------------------------------

\subsection{Kernel specification}

All simulations use:
\begin{itemize}
  \item Kernel: \texttt{cog\_v3\_octavian240\_multiplicative\_v1}
  \item Convention: \texttt{v3\_octavian240\_perm\_0-1-2-3-4-7-5-6\_sign\_all\_plus}
  \item Stencil: \texttt{cube26} (26-neighbour cubic lattice)
  \item Event order: \texttt{synchronous\_parallel\_v1}
\end{itemize}

The state space is $S_{960} = C_4 \times Q_{240}$ for Phase~D runs (C12 integration
in progress), with $S_{2880} = C_{12} \times Q_{240}$ probes for the R3 measurements.

\subsection{R3 staircase: the three probes}

\textit{Mode 1 — Matched (broken):} The kick and state elements share the same $C_{12}$
phase $p$, so the seed phase doubles to $2p \bmod 12$ (always even).  The $\Delta p = 3$
channel is algebraically unreachable from even phases.  Result: $R_3 = 0.000$.

\textit{Mode 2 — Zero-kick fix:} The kick is fixed at phase 0 ($\texttt{kick\_id} = \texttt{\_s\_id}(0, q_k, 240)$),
allowing odd-phase seeds.  The $\Delta p = 3$ channel becomes accessible.
Result: $R_3 = 0.090$, $A_1 = -0.036$.

\textit{Mode 3 — Zero-kick + odd lane:} Seed phases are additionally restricted to
$p \equiv 1 \pmod 4$ (the odd-lane, maximally activating $\Delta p = 3$).
Result: $R_3 = 0.220$, $A_1 = -0.046$.

\subsection{Two-region d3 probe}

On a $16 \times 16 \times 16$ grid with half the cells seeded at phase $p=0$ and half
at $p=3$, the $\texttt{axial6}$ stencil yields:
\begin{equation}
  N_{d3}(\text{region 1}) = 38{,}919, \quad N_{d3}(\text{region 2}) = 37{,}282 .
\end{equation}
Both regions show d3 hops, confirming that the $\Delta p = 3$ channel is \emph{physically
available} (not forbidden) and that $R_3 = 0$ in Mode~1 is due to the seeding bug,
not an algebraic impossibility.

\subsection{Period-48 photon motif}

After 576 overnight batches with the $Q_{240}$-only multiplicative kernel:
\begin{center}
\begin{tabular}{ll}
  \toprule
  Parameter & Value \\
  \midrule
  Period $N$ & 48 \\
  Seed family & \texttt{sheet\_y} \\
  Kick element & \texttt{1*e011} (= $e_3$ in octonion basis) \\
  Score & 0.5656 \\
  Gate status & gate0$\checkmark$, gate1$\checkmark$, gate3$\checkmark$, gate2$\times$, gate4$\times$ \\
  Gate2 reason & $\mathrm{both\_hit\_fraction} = 1.0$ (wave-like, expected for photon) \\
  Gate4 reason & $A_\chi = 0.0$ (self-conjugate boson, expected for photon) \\
  \bottomrule
\end{tabular}
\end{center}

Under the proposed boson gate criteria (gate2\_boson: $\mathrm{both\_hit} > 0.9$;
gate4\_boson: $|A_\chi| < 0.05$), this candidate passes all relevant gates.

%-------------------------------------------------------
\section{PSL(2,7) Orbit Data}
\label{app:psl27}
%-------------------------------------------------------

From the order-6 conjugation probe (RFC-013):
\begin{itemize}
  \item Total order-6 elements: 168 ✓
  \item Tested family: 56-element order-3 conjugation class
  \item Orbit partition: $168 = 56 + 56 + 56$ (three orbits of 56)
  \item Stabiliser: every element fixed by exactly 2 of the 56 conjugators
  \item Closure: $\mathrm{closure\_ok} = \mathrm{False}$ (expected: the 56 tested do not
        close under multiplication)
  \item Faithful: $\mathrm{faithful\_ok} = \mathrm{False}$ (expected: 56-element family is
        a subclass, not all of PSL(2,7))
\end{itemize}

Interpretation: the 56-element tested family is the \emph{order-3 conjugacy class} of
$\mathrm{PSL}(2,7)$.  Acting with only this class partitions the 168 order-6 elements
into three cosets corresponding to the three generation sectors (Gen1, Gen2, Gen3).
The full PSL(2,7) action (Family~A test, pending) is predicted to give a single orbit
of 168 with trivial stabilisers.

%-------------------------------------------------------
\section{Computational Verification: Python Code}
\label{app:code}
%-------------------------------------------------------

The following results are verified by \texttt{calc/derive\_sm\_quantum\_numbers.py}
using Python \texttt{Fraction} arithmetic (exact rational arithmetic):

\begin{verbatim}
from fractions import Fraction

N_c = 3; N_w = 2
# Gauge boson count
N_SU3 = N_c**2 - 1       # = 8
N_SU2 = N_w**2 - 1       # = 3
N_gauge = N_SU3 + N_SU2 + 1  # = 12  ✓

# Weinberg angle
k = Fraction(N_c, N_c + N_w)   # = 3/5
sin2_theta_W = k / (1 + k)     # = 3/8  ✓
assert sin2_theta_W == Fraction(3, 8)

# QCD beta function
N_f = 3
beta0_UV = 11 - Fraction(2, 3) * N_f  # = 9  ✓
\end{verbatim}

All assertions pass. Full test suite: \texttt{pytest calc/test\_derive\_sm\_quantum\_numbers.py}
(24 tests, all passing).

% Encoding: ASCII-safe LaTeX

\appendix

%-------------------------------------------------------
\section{Why $C_{12}$: Minimality Argument}
\label{app:why_c12}
%-------------------------------------------------------

\begin{proposition}[Minimality of $C_{12}$]
  $C_{12}$ is the smallest cyclic group $C_n$ that contains both:
  (i) a $\mathbb{Z}_4$ subgroup (required for $\mathrm{U}(1)_Y$ hypercharge structure),
  (ii) a $\mathbb{Z}_3$ subgroup (required for three fermion generations).
\end{proposition}

\begin{proof}
  A cyclic group $C_n = \mathbb{Z}_n$ contains a $\mathbb{Z}_k$ subgroup if and only if
  $k \mid n$.  We require both $4 \mid n$ and $3 \mid n$, hence $12 = \mathrm{lcm}(4,3) \mid n$.
  The smallest such $n$ is $n = 12$.  Furthermore, when $n = 12$, by the Chinese Remainder
  Theorem $\mathbb{Z}_{12} \cong \mathbb{Z}_4 \times \mathbb{Z}_3$ (since $\gcd(4,3)=1$),
  giving exact product structure.  For $n = 24$: $\mathbb{Z}_{24} \cong \mathbb{Z}_8 \times \mathbb{Z}_3$,
  which has $\mathbb{Z}_3$ but only $\mathbb{Z}_4 \subset \mathbb{Z}_8$ (no separate $\mathbb{Z}_4$ factor).
  Therefore $C_{12}$ is the unique minimal choice with $\mathbb{Z}_{12} \cong \mathbb{Z}_4 \times \mathbb{Z}_3$
  as a direct product.
\end{proof}

\begin{remark}
  The $\mathbb{Z}_4$ factor requires explanation.  In the old $S_{960} = C_4 \times Q_{240}$
  model~\cite{RFC009_S960}, the 4-element clock encoded the hypercharge quantum number.
  The states $\{1, i, -1, -i\}$ correspond to hypercharge $Y \in \{0, +1, 0, -1\}$ for
  the four complex phases.  The $C_{12}$ extension preserves this $\mathbb{Z}_4$ sub-structure
  (each generation has phases $\{p, p+3, p+6, p+9\}$ forming a $\mathbb{Z}_4$ orbit at fixed
  generation $g = p \bmod 3$), while adding the $\mathbb{Z}_3$ generation sector.

  The 12th cyclotomic field $\mathbb{Q}(\zeta_{12}) = \mathbb{Q}(i, \sqrt{3})$ is the minimal
  number field containing both the $\mathrm{SU}(2)$ eigenvalue $i$ (from $\mathbb{Z}_4$) and
  the $\mathrm{SU}(3)$ eigenvalue $\omega = e^{2\pi i/3}$ (from $\mathbb{Z}_3$).  The C$_{12}$
  roots of unity span both sub-structures simultaneously.
\end{remark}

%-------------------------------------------------------
\section{Explicit Weinberg Angle Derivation}
\label{app:weinberg}
%-------------------------------------------------------

We provide a self-contained derivation of $\sin^2\!\theta_W = 3/8$ that does not invoke
the full SU(5) GUT machinery, working directly from the $S_{2880}$ quantum number structure.

\subsection{Setup}

In $S_{2880}$, the electric charge operator is
\begin{equation}
  Q = T_3 + \tfrac{1}{2} Y ,
\label{eq:charge_def}
\end{equation}
where $T_3$ is the third component of weak isospin and $Y$ is the hypercharge.

For the left-handed lepton doublet $(\nu_L, e_L)$ at generation $g$ in $S_{2880}$:
\begin{equation}
  T_3(\nu_L) = +\tfrac{1}{2}, \quad T_3(e_L) = -\tfrac{1}{2}, \quad Y(\nu_L) = Y(e_L) = -1 .
\end{equation}
This gives $Q(\nu_L) = 0$, $Q(e_L) = -1$ as required.

\subsection{The mixing angle from normalisation}

The $\mathrm{U}(1)_Y$ coupling $g'$ and $\mathrm{SU}(2)_L$ coupling $g$ are related to
the electromagnetic coupling $e$ by
\begin{equation}
  e = g \sin\theta_W = g' \cos\theta_W .
\end{equation}

In the $S_{2880}$ model, the symmetry group is embedded in $G_2 \supset \mathrm{SU}(3)_c
\times \mathrm{SU}(2)_L$.  The $\mathrm{U}(1)_Y$ hypercharge generator is the \emph{trace}
direction orthogonal to $\mathrm{SU}(2)_L$ in the embedding.

\subsection{Derivation via gauge boson counting}

The Weinberg angle satisfies (at GUT scale, where $g = g'$ in the normalised sense):
\begin{align}
  \sin^2\!\theta_W &= \frac{g'^2}{g^2 + g'^2}
  = \frac{\mathrm{dim}(\mathrm{U}(1)_Y)}{\mathrm{dim}(\mathrm{U}(1)_Y) + \mathrm{dim}(\mathrm{SU}(2)_L)}
  = \frac{1}{1 + N_w^2 - 1} = \frac{1}{N_w^2} .
\end{align}
But $N_w = 2$, giving $\sin^2\!\theta_W = 1/4$ only (the standard SU(5) naive value).

The correct normalisation in SU(5) uses the Dynkin index of the $\mathrm{U}(1)_Y$
generator within the adjoint of SU(5).  The standard result is:
\begin{equation}
  \sin^2\!\theta_W = \frac{g'^2}{g^2 + g'^2}
  \Bigg|_{\mathrm{SU}(5)\,\mathrm{normalised}}
  = \frac{3/5}{1 + 3/5} = \frac{3}{8} .
\end{equation}
The factor $3/5$ arises because the hypercharge generator $Y$ in SU(5) has normalisation
$k = N_c/(N_c + N_w) = 3/5$.  Substituting:
\begin{equation}
  \sin^2\!\theta_W = \frac{k}{1+k} = \frac{3/5}{1+3/5} = \frac{3/5}{8/5} = \frac{3}{8} .
\end{equation}

\subsection{Alternative derivation from charge sum rules}

Equivalently, for one complete generation of 15 left-handed Weyl fermions in the SM
(2 leptons + $3 \times 3$ quarks + right-handed partners), the anomaly cancellation
conditions uniquely fix $\sin^2\theta_W = 3/8$ at the GUT scale.
In $S_{2880}$, the 15 Weyl fermions per generation are accounted for by:
\begin{itemize}
  \item 2 lepton states $(\nu_L, e_L)$ per generation
  \item $3 \times 3 = 9$ quark states ($u_L, d_L$ doublet $\times$ 3 colours, plus $u_R, d_R$)
  \item $\mathrm{Tr}[Y] = 0$, $\mathrm{Tr}[Y^3] = 0$ enforced by $Q_{240}$ Witt structure
\end{itemize}

Verification: \texttt{derive\_sm\_quantum\_numbers.py::derive\_anomaly\_cancellation()} confirms
$\mathrm{Tr}[Y] = 0$ and $\mathrm{Tr}[Y^3] = 0$ for the 15 LH fermions using Python
\texttt{Fraction} arithmetic (no floating-point error).

%-------------------------------------------------------
\section{Koide Formula: Full Calculation}
\label{app:koide_full}
%-------------------------------------------------------

We provide the explicit calculation confirming Theorem~\ref{thm:Koide} with all
intermediate steps, for the benefit of readers unfamiliar with the Brannen
parameterisation.

\subsection{The three angles}

For $k = 0, 1, 2$: $\phi_k = 2\pi k / 3 + \delta$.

The three angles are $\phi_0 = \delta$, $\phi_1 = 2\pi/3 + \delta$, $\phi_2 = 4\pi/3 + \delta$.

\subsection{Sum of square roots}

\begin{align}
  S_1 &= \sum_{k=0}^2 \sqrt{m_k} = C \sum_{k=0}^2 (1 + \sqrt{2} \cos\phi_k) \notag \\
      &= C \left(3 + \sqrt{2} \sum_{k=0}^2 \cos\phi_k\right) .
\end{align}

\begin{lemma}
  For any $\delta$: $\sum_{k=0}^2 \cos(2\pi k/3 + \delta) = 0$.
\end{lemma}

\begin{proof}
  $\sum_{k=0}^2 \cos(2\pi k/3 + \delta) = \mathrm{Re}\left(e^{i\delta} \sum_{k=0}^2 e^{2\pi i k/3}\right) = 0$,
  since $\sum_{k=0}^2 \omega^k = 0$ for $\omega = e^{2\pi i/3} \neq 1$.
\end{proof}

Therefore $S_1 = 3C$.

\subsection{Sum of masses}

\begin{align}
  S_2 &= \sum_{k=0}^2 m_k = C^2 \sum_{k=0}^2 (1 + \sqrt{2}\cos\phi_k)^2 \notag \\
      &= C^2 \sum_{k=0}^2 \left(1 + 2\sqrt{2}\cos\phi_k + 2\cos^2\!\phi_k\right) \notag \\
      &= C^2 \left(3 + 2\sqrt{2} \cdot 0 + 2 \sum_{k=0}^2 \cos^2\!\phi_k\right) .
\end{align}

\begin{lemma}
  For any $\delta$: $\sum_{k=0}^2 \cos^2(2\pi k/3 + \delta) = 3/2$.
\end{lemma}

\begin{proof}
  Using $\cos^2\theta = (1 + \cos 2\theta)/2$:
  $\sum_{k=0}^2 \cos^2\phi_k = \frac{1}{2}\left(3 + \sum_{k=0}^2 \cos(4\pi k/3 + 2\delta)\right) = \frac{3}{2} + 0 = \frac{3}{2}$,
  since $\sum_{k=0}^2 \cos(4\pi k/3 + 2\delta) = 0$ by the same argument as above
  (with $\omega = e^{4\pi i/3} = e^{-2\pi i/3}$ also a primitive cube root of unity).
\end{proof}

Therefore $S_2 = C^2(3 + 3) = 6C^2$.

\subsection{Koide ratio}

\begin{equation}
  K = \frac{S_2}{S_1^2} = \frac{6C^2}{9C^2} = \frac{2}{3} . \qed
\end{equation}

\subsection{Physical check for charged leptons}

Experimental lepton masses (PDG 2024):
$m_e = 0.511$~MeV, $m_\mu = 105.66$~MeV, $m_\tau = 1776.86$~MeV.
\begin{equation}
  K_{\mathrm{exp}} = \frac{m_e + m_\mu + m_\tau}{(\sqrt{m_e} + \sqrt{m_\mu} + \sqrt{m_\tau})^2}
  = \frac{1883.03}{(0.7149 + 10.279 + 42.153)^2} = \frac{1883.03}{2824.1} \approx 0.6668 \approx 2/3 .
\end{equation}
Agreement to five significant figures.

%-------------------------------------------------------
\section{Phase~D Simulation Details}
\label{app:simulation}
%-------------------------------------------------------

\subsection{Kernel specification}

All simulations use:
\begin{itemize}
  \item Kernel: \texttt{cog\_v3\_octavian240\_multiplicative\_v1}
  \item Convention: \texttt{v3\_octavian240\_perm\_0-1-2-3-4-7-5-6\_sign\_all\_plus}
  \item Stencil: \texttt{cube26} (26-neighbour cubic lattice)
  \item Event order: \texttt{synchronous\_parallel\_v1}
\end{itemize}

The state space is $S_{960} = C_4 \times Q_{240}$ for Phase~D runs (C12 integration
in progress), with $S_{2880} = C_{12} \times Q_{240}$ probes for the R3 measurements.

\subsection{R3 staircase: the three probes}

\textit{Mode 1 — Matched (broken):} The kick and state elements share the same $C_{12}$
phase $p$, so the seed phase doubles to $2p \bmod 12$ (always even).  The $\Delta p = 3$
channel is algebraically unreachable from even phases.  Result: $R_3 = 0.000$.

\textit{Mode 2 — Zero-kick fix:} The kick is fixed at phase 0 ($\texttt{kick\_id} = \texttt{\_s\_id}(0, q_k, 240)$),
allowing odd-phase seeds.  The $\Delta p = 3$ channel becomes accessible.
Result: $R_3 = 0.090$, $A_1 = -0.036$.

\textit{Mode 3 — Zero-kick + odd lane:} Seed phases are additionally restricted to
$p \equiv 1 \pmod 4$ (the odd-lane, maximally activating $\Delta p = 3$).
Result: $R_3 = 0.220$, $A_1 = -0.046$.

\subsection{Two-region d3 probe}

On a $16 \times 16 \times 16$ grid with half the cells seeded at phase $p=0$ and half
at $p=3$, the $\texttt{axial6}$ stencil yields:
\begin{equation}
  N_{d3}(\text{region 1}) = 38{,}919, \quad N_{d3}(\text{region 2}) = 37{,}282 .
\end{equation}
Both regions show d3 hops, confirming that the $\Delta p = 3$ channel is \emph{physically
available} (not forbidden) and that $R_3 = 0$ in Mode~1 is due to the seeding bug,
not an algebraic impossibility.

\subsection{Period-48 photon motif}

After 576 overnight batches with the $Q_{240}$-only multiplicative kernel:
\begin{center}
\begin{tabular}{ll}
  \toprule
  Parameter & Value \\
  \midrule
  Period $N$ & 48 \\
  Seed family & \texttt{sheet\_y} \\
  Kick element & \texttt{1*e011} (= $e_3$ in octonion basis) \\
  Score & 0.5656 \\
  Gate status & gate0$\checkmark$, gate1$\checkmark$, gate3$\checkmark$, gate2$\times$, gate4$\times$ \\
  Gate2 reason & $\mathrm{both\_hit\_fraction} = 1.0$ (wave-like, expected for photon) \\
  Gate4 reason & $A_\chi = 0.0$ (self-conjugate boson, expected for photon) \\
  \bottomrule
\end{tabular}
\end{center}

Under the proposed boson gate criteria (gate2\_boson: $\mathrm{both\_hit} > 0.9$;
gate4\_boson: $|A_\chi| < 0.05$), this candidate passes all relevant gates.

%-------------------------------------------------------
\section{PSL(2,7) Orbit Data}
\label{app:psl27}
%-------------------------------------------------------

From the order-6 conjugation probe (RFC-013):
\begin{itemize}
  \item Total order-6 elements: 168 ✓
  \item Tested family: 56-element order-3 conjugation class
  \item Orbit partition: $168 = 56 + 56 + 56$ (three orbits of 56)
  \item Stabiliser: every element fixed by exactly 2 of the 56 conjugators
  \item Closure: $\mathrm{closure\_ok} = \mathrm{False}$ (expected: the 56 tested do not
        close under multiplication)
  \item Faithful: $\mathrm{faithful\_ok} = \mathrm{False}$ (expected: 56-element family is
        a subclass, not all of PSL(2,7))
\end{itemize}

Interpretation: the 56-element tested family is the \emph{order-3 conjugacy class} of
$\mathrm{PSL}(2,7)$.  Acting with only this class partitions the 168 order-6 elements
into three cosets corresponding to the three generation sectors (Gen1, Gen2, Gen3).
The full PSL(2,7) action (Family~A test, pending) is predicted to give a single orbit
of 168 with trivial stabilisers.

%-------------------------------------------------------
\section{Computational Verification: Python Code}
\label{app:code}
%-------------------------------------------------------

The following results are verified by \texttt{calc/derive\_sm\_quantum\_numbers.py}
using Python \texttt{Fraction} arithmetic (exact rational arithmetic):

\begin{verbatim}
from fractions import Fraction

N_c = 3; N_w = 2
# Gauge boson count
N_SU3 = N_c**2 - 1       # = 8
N_SU2 = N_w**2 - 1       # = 3
N_gauge = N_SU3 + N_SU2 + 1  # = 12  ✓

# Weinberg angle
k = Fraction(N_c, N_c + N_w)   # = 3/5
sin2_theta_W = k / (1 + k)     # = 3/8  ✓
assert sin2_theta_W == Fraction(3, 8)

# QCD beta function
N_f = 3
beta0_UV = 11 - Fraction(2, 3) * N_f  # = 9  ✓
\end{verbatim}

All assertions pass. Full test suite: \texttt{pytest calc/test\_derive\_sm\_quantum\_numbers.py}
(24 tests, all passing).

% Encoding: ASCII-safe LaTeX

\appendix

%-------------------------------------------------------
\section{Why $C_{12}$: Minimality Argument}
\label{app:why_c12}
%-------------------------------------------------------

\begin{proposition}[Minimality of $C_{12}$]
  $C_{12}$ is the smallest cyclic group $C_n$ that contains both:
  (i) a $\mathbb{Z}_4$ subgroup (required for $\mathrm{U}(1)_Y$ hypercharge structure),
  (ii) a $\mathbb{Z}_3$ subgroup (required for three fermion generations).
\end{proposition}

\begin{proof}
  A cyclic group $C_n = \mathbb{Z}_n$ contains a $\mathbb{Z}_k$ subgroup if and only if
  $k \mid n$.  We require both $4 \mid n$ and $3 \mid n$, hence $12 = \mathrm{lcm}(4,3) \mid n$.
  The smallest such $n$ is $n = 12$.  Furthermore, when $n = 12$, by the Chinese Remainder
  Theorem $\mathbb{Z}_{12} \cong \mathbb{Z}_4 \times \mathbb{Z}_3$ (since $\gcd(4,3)=1$),
  giving exact product structure.  For $n = 24$: $\mathbb{Z}_{24} \cong \mathbb{Z}_8 \times \mathbb{Z}_3$,
  which has $\mathbb{Z}_3$ but only $\mathbb{Z}_4 \subset \mathbb{Z}_8$ (no separate $\mathbb{Z}_4$ factor).
  Therefore $C_{12}$ is the unique minimal choice with $\mathbb{Z}_{12} \cong \mathbb{Z}_4 \times \mathbb{Z}_3$
  as a direct product.
\end{proof}

\begin{remark}
  The $\mathbb{Z}_4$ factor requires explanation.  In the old $S_{960} = C_4 \times Q_{240}$
  model~\cite{RFC009_S960}, the 4-element clock encoded the hypercharge quantum number.
  The states $\{1, i, -1, -i\}$ correspond to hypercharge $Y \in \{0, +1, 0, -1\}$ for
  the four complex phases.  The $C_{12}$ extension preserves this $\mathbb{Z}_4$ sub-structure
  (each generation has phases $\{p, p+3, p+6, p+9\}$ forming a $\mathbb{Z}_4$ orbit at fixed
  generation $g = p \bmod 3$), while adding the $\mathbb{Z}_3$ generation sector.

  The 12th cyclotomic field $\mathbb{Q}(\zeta_{12}) = \mathbb{Q}(i, \sqrt{3})$ is the minimal
  number field containing both the $\mathrm{SU}(2)$ eigenvalue $i$ (from $\mathbb{Z}_4$) and
  the $\mathrm{SU}(3)$ eigenvalue $\omega = e^{2\pi i/3}$ (from $\mathbb{Z}_3$).  The C$_{12}$
  roots of unity span both sub-structures simultaneously.
\end{remark}

%-------------------------------------------------------
\section{Explicit Weinberg Angle Derivation}
\label{app:weinberg}
%-------------------------------------------------------

We provide a self-contained derivation of $\sin^2\!\theta_W = 3/8$ that does not invoke
the full SU(5) GUT machinery, working directly from the $S_{2880}$ quantum number structure.

\subsection{Setup}

In $S_{2880}$, the electric charge operator is
\begin{equation}
  Q = T_3 + \tfrac{1}{2} Y ,
\label{eq:charge_def}
\end{equation}
where $T_3$ is the third component of weak isospin and $Y$ is the hypercharge.

For the left-handed lepton doublet $(\nu_L, e_L)$ at generation $g$ in $S_{2880}$:
\begin{equation}
  T_3(\nu_L) = +\tfrac{1}{2}, \quad T_3(e_L) = -\tfrac{1}{2}, \quad Y(\nu_L) = Y(e_L) = -1 .
\end{equation}
This gives $Q(\nu_L) = 0$, $Q(e_L) = -1$ as required.

\subsection{The mixing angle from normalisation}

The $\mathrm{U}(1)_Y$ coupling $g'$ and $\mathrm{SU}(2)_L$ coupling $g$ are related to
the electromagnetic coupling $e$ by
\begin{equation}
  e = g \sin\theta_W = g' \cos\theta_W .
\end{equation}

In the $S_{2880}$ model, the symmetry group is embedded in $G_2 \supset \mathrm{SU}(3)_c
\times \mathrm{SU}(2)_L$.  The $\mathrm{U}(1)_Y$ hypercharge generator is the \emph{trace}
direction orthogonal to $\mathrm{SU}(2)_L$ in the embedding.

\subsection{Derivation via gauge boson counting}

The Weinberg angle satisfies (at GUT scale, where $g = g'$ in the normalised sense):
\begin{align}
  \sin^2\!\theta_W &= \frac{g'^2}{g^2 + g'^2}
  = \frac{\mathrm{dim}(\mathrm{U}(1)_Y)}{\mathrm{dim}(\mathrm{U}(1)_Y) + \mathrm{dim}(\mathrm{SU}(2)_L)}
  = \frac{1}{1 + N_w^2 - 1} = \frac{1}{N_w^2} .
\end{align}
But $N_w = 2$, giving $\sin^2\!\theta_W = 1/4$ only (the standard SU(5) naive value).

The correct normalisation in SU(5) uses the Dynkin index of the $\mathrm{U}(1)_Y$
generator within the adjoint of SU(5).  The standard result is:
\begin{equation}
  \sin^2\!\theta_W = \frac{g'^2}{g^2 + g'^2}
  \Bigg|_{\mathrm{SU}(5)\,\mathrm{normalised}}
  = \frac{3/5}{1 + 3/5} = \frac{3}{8} .
\end{equation}
The factor $3/5$ arises because the hypercharge generator $Y$ in SU(5) has normalisation
$k = N_c/(N_c + N_w) = 3/5$.  Substituting:
\begin{equation}
  \sin^2\!\theta_W = \frac{k}{1+k} = \frac{3/5}{1+3/5} = \frac{3/5}{8/5} = \frac{3}{8} .
\end{equation}

\subsection{Alternative derivation from charge sum rules}

Equivalently, for one complete generation of 15 left-handed Weyl fermions in the SM
(2 leptons + $3 \times 3$ quarks + right-handed partners), the anomaly cancellation
conditions uniquely fix $\sin^2\theta_W = 3/8$ at the GUT scale.
In $S_{2880}$, the 15 Weyl fermions per generation are accounted for by:
\begin{itemize}
  \item 2 lepton states $(\nu_L, e_L)$ per generation
  \item $3 \times 3 = 9$ quark states ($u_L, d_L$ doublet $\times$ 3 colours, plus $u_R, d_R$)
  \item $\mathrm{Tr}[Y] = 0$, $\mathrm{Tr}[Y^3] = 0$ enforced by $Q_{240}$ Witt structure
\end{itemize}

Verification: \texttt{derive\_sm\_quantum\_numbers.py::derive\_anomaly\_cancellation()} confirms
$\mathrm{Tr}[Y] = 0$ and $\mathrm{Tr}[Y^3] = 0$ for the 15 LH fermions using Python
\texttt{Fraction} arithmetic (no floating-point error).

%-------------------------------------------------------
\section{Koide Formula: Full Calculation}
\label{app:koide_full}
%-------------------------------------------------------

We provide the explicit calculation confirming Theorem~\ref{thm:Koide} with all
intermediate steps, for the benefit of readers unfamiliar with the Brannen
parameterisation.

\subsection{The three angles}

For $k = 0, 1, 2$: $\phi_k = 2\pi k / 3 + \delta$.

The three angles are $\phi_0 = \delta$, $\phi_1 = 2\pi/3 + \delta$, $\phi_2 = 4\pi/3 + \delta$.

\subsection{Sum of square roots}

\begin{align}
  S_1 &= \sum_{k=0}^2 \sqrt{m_k} = C \sum_{k=0}^2 (1 + \sqrt{2} \cos\phi_k) \notag \\
      &= C \left(3 + \sqrt{2} \sum_{k=0}^2 \cos\phi_k\right) .
\end{align}

\begin{lemma}
  For any $\delta$: $\sum_{k=0}^2 \cos(2\pi k/3 + \delta) = 0$.
\end{lemma}

\begin{proof}
  $\sum_{k=0}^2 \cos(2\pi k/3 + \delta) = \mathrm{Re}\left(e^{i\delta} \sum_{k=0}^2 e^{2\pi i k/3}\right) = 0$,
  since $\sum_{k=0}^2 \omega^k = 0$ for $\omega = e^{2\pi i/3} \neq 1$.
\end{proof}

Therefore $S_1 = 3C$.

\subsection{Sum of masses}

\begin{align}
  S_2 &= \sum_{k=0}^2 m_k = C^2 \sum_{k=0}^2 (1 + \sqrt{2}\cos\phi_k)^2 \notag \\
      &= C^2 \sum_{k=0}^2 \left(1 + 2\sqrt{2}\cos\phi_k + 2\cos^2\!\phi_k\right) \notag \\
      &= C^2 \left(3 + 2\sqrt{2} \cdot 0 + 2 \sum_{k=0}^2 \cos^2\!\phi_k\right) .
\end{align}

\begin{lemma}
  For any $\delta$: $\sum_{k=0}^2 \cos^2(2\pi k/3 + \delta) = 3/2$.
\end{lemma}

\begin{proof}
  Using $\cos^2\theta = (1 + \cos 2\theta)/2$:
  $\sum_{k=0}^2 \cos^2\phi_k = \frac{1}{2}\left(3 + \sum_{k=0}^2 \cos(4\pi k/3 + 2\delta)\right) = \frac{3}{2} + 0 = \frac{3}{2}$,
  since $\sum_{k=0}^2 \cos(4\pi k/3 + 2\delta) = 0$ by the same argument as above
  (with $\omega = e^{4\pi i/3} = e^{-2\pi i/3}$ also a primitive cube root of unity).
\end{proof}

Therefore $S_2 = C^2(3 + 3) = 6C^2$.

\subsection{Koide ratio}

\begin{equation}
  K = \frac{S_2}{S_1^2} = \frac{6C^2}{9C^2} = \frac{2}{3} . \qed
\end{equation}

\subsection{Physical check for charged leptons}

Experimental lepton masses (PDG 2024):
$m_e = 0.511$~MeV, $m_\mu = 105.66$~MeV, $m_\tau = 1776.86$~MeV.
\begin{equation}
  K_{\mathrm{exp}} = \frac{m_e + m_\mu + m_\tau}{(\sqrt{m_e} + \sqrt{m_\mu} + \sqrt{m_\tau})^2}
  = \frac{1883.03}{(0.7149 + 10.279 + 42.153)^2} = \frac{1883.03}{2824.1} \approx 0.6668 \approx 2/3 .
\end{equation}
Agreement to five significant figures.

%-------------------------------------------------------
\section{Phase~D Simulation Details}
\label{app:simulation}
%-------------------------------------------------------

\subsection{Kernel specification}

All simulations use:
\begin{itemize}
  \item Kernel: \texttt{cog\_v3\_octavian240\_multiplicative\_v1}
  \item Convention: \texttt{v3\_octavian240\_perm\_0-1-2-3-4-7-5-6\_sign\_all\_plus}
  \item Stencil: \texttt{cube26} (26-neighbour cubic lattice)
  \item Event order: \texttt{synchronous\_parallel\_v1}
\end{itemize}

The state space is $S_{960} = C_4 \times Q_{240}$ for Phase~D runs (C12 integration
in progress), with $S_{2880} = C_{12} \times Q_{240}$ probes for the R3 measurements.

\subsection{R3 staircase: the three probes}

\textit{Mode 1 — Matched (broken):} The kick and state elements share the same $C_{12}$
phase $p$, so the seed phase doubles to $2p \bmod 12$ (always even).  The $\Delta p = 3$
channel is algebraically unreachable from even phases.  Result: $R_3 = 0.000$.

\textit{Mode 2 — Zero-kick fix:} The kick is fixed at phase 0 ($\texttt{kick\_id} = \texttt{\_s\_id}(0, q_k, 240)$),
allowing odd-phase seeds.  The $\Delta p = 3$ channel becomes accessible.
Result: $R_3 = 0.090$, $A_1 = -0.036$.

\textit{Mode 3 — Zero-kick + odd lane:} Seed phases are additionally restricted to
$p \equiv 1 \pmod 4$ (the odd-lane, maximally activating $\Delta p = 3$).
Result: $R_3 = 0.220$, $A_1 = -0.046$.

\subsection{Two-region d3 probe}

On a $16 \times 16 \times 16$ grid with half the cells seeded at phase $p=0$ and half
at $p=3$, the $\texttt{axial6}$ stencil yields:
\begin{equation}
  N_{d3}(\text{region 1}) = 38{,}919, \quad N_{d3}(\text{region 2}) = 37{,}282 .
\end{equation}
Both regions show d3 hops, confirming that the $\Delta p = 3$ channel is \emph{physically
available} (not forbidden) and that $R_3 = 0$ in Mode~1 is due to the seeding bug,
not an algebraic impossibility.

\subsection{Period-48 photon motif}

After 576 overnight batches with the $Q_{240}$-only multiplicative kernel:
\begin{center}
\begin{tabular}{ll}
  \toprule
  Parameter & Value \\
  \midrule
  Period $N$ & 48 \\
  Seed family & \texttt{sheet\_y} \\
  Kick element & \texttt{1*e011} (= $e_3$ in octonion basis) \\
  Score & 0.5656 \\
  Gate status & gate0$\checkmark$, gate1$\checkmark$, gate3$\checkmark$, gate2$\times$, gate4$\times$ \\
  Gate2 reason & $\mathrm{both\_hit\_fraction} = 1.0$ (wave-like, expected for photon) \\
  Gate4 reason & $A_\chi = 0.0$ (self-conjugate boson, expected for photon) \\
  \bottomrule
\end{tabular}
\end{center}

Under the proposed boson gate criteria (gate2\_boson: $\mathrm{both\_hit} > 0.9$;
gate4\_boson: $|A_\chi| < 0.05$), this candidate passes all relevant gates.

%-------------------------------------------------------
\section{PSL(2,7) Orbit Data}
\label{app:psl27}
%-------------------------------------------------------

From the order-6 conjugation probe (RFC-013):
\begin{itemize}
  \item Total order-6 elements: 168 ✓
  \item Tested family: 56-element order-3 conjugation class
  \item Orbit partition: $168 = 56 + 56 + 56$ (three orbits of 56)
  \item Stabiliser: every element fixed by exactly 2 of the 56 conjugators
  \item Closure: $\mathrm{closure\_ok} = \mathrm{False}$ (expected: the 56 tested do not
        close under multiplication)
  \item Faithful: $\mathrm{faithful\_ok} = \mathrm{False}$ (expected: 56-element family is
        a subclass, not all of PSL(2,7))
\end{itemize}

Interpretation: the 56-element tested family is the \emph{order-3 conjugacy class} of
$\mathrm{PSL}(2,7)$.  Acting with only this class partitions the 168 order-6 elements
into three cosets corresponding to the three generation sectors (Gen1, Gen2, Gen3).
The full PSL(2,7) action (Family~A test, pending) is predicted to give a single orbit
of 168 with trivial stabilisers.

%-------------------------------------------------------
\section{Computational Verification: Python Code}
\label{app:code}
%-------------------------------------------------------

The following results are verified by \texttt{calc/derive\_sm\_quantum\_numbers.py}
using Python \texttt{Fraction} arithmetic (exact rational arithmetic):

\begin{verbatim}
from fractions import Fraction

N_c = 3; N_w = 2
# Gauge boson count
N_SU3 = N_c**2 - 1       # = 8
N_SU2 = N_w**2 - 1       # = 3
N_gauge = N_SU3 + N_SU2 + 1  # = 12  ✓

# Weinberg angle
k = Fraction(N_c, N_c + N_w)   # = 3/5
sin2_theta_W = k / (1 + k)     # = 3/8  ✓
assert sin2_theta_W == Fraction(3, 8)

# QCD beta function
N_f = 3
beta0_UV = 11 - Fraction(2, 3) * N_f  # = 9  ✓
\end{verbatim}

All assertions pass. Full test suite: \texttt{pytest calc/test\_derive\_sm\_quantum\_numbers.py}
(24 tests, all passing).
